\begin{result}{E.2.30}{The embedded completion determined by all proper cover lattices}{Proved}
Let \(F=\langle x,y\rangle\) be the specified topological fundamental
group of \(X=\mathbf P^1(\C)\setminus\{0,1,\infty\}\), put
\(z=(xy)^{-1}\), and put \(G=\widehat F\). For every open normal
subgroup \(N\subset G\), let \(Y_N\) be the actual regular cover,
let \(C_N\) be its smooth compactification, and define
\[
\begin{aligned}
 K_N&=F\cap N,& A_N^c&=H_1(C_N,\Z),\\
 L_N&=\ker(K_N\longrightarrow A_N^c),&
 J_N&=\overline{L_N}^{\,G},\\
 \Gamma_N^c&=F/L_N,&P_N^c&=G/J_N,& Q_N&=G/N.
\end{aligned}
\tag{E.2.30a}
\]
There are actual embeddings and inverse-limit identifications
\[
 \Gamma_N^c\hookrightarrow P_N^c=\widehat{\Gamma_N^c},
 \qquad
 \mathcal H_c:=\bigcap_N(\pi_N^c)^{-1}(\Gamma_N^c)
       \simeq\varprojlim_N\Gamma_N^c\subset G,
 \quad \pi_N^c:G\longrightarrow P_N^c.
\tag{E.2.30b}
\]
The subgroup \(\mathcal H_c\) is self-normalizing in \(G\).

Let \(\alpha\) be an actual continuous automorphism of \(G\)
that permutes the three peripheral conjugacy classes and sends
each specified peripheral generator to a conjugate of a unit
power of the corresponding generator. For an open subgroup
\(U\subset G\), write \(C_U\) for the compactification of its
actual cover and \(\Lambda_U^c=H_1(C_U,\Z)\). The following
conditions are equivalent:
\begin{enumerate}
\item The actual induced completed proper-homology map carries
\(\Lambda_U^c\) onto \(\Lambda_{\alpha(U)}^c\) for every open
\(U\subset G\).
\item It carries \(A_N^c\) onto itself on one characteristic
cofinal chain of open normal subgroups.
\item There is one \(h\in G\) with
\((\operatorname{ad}(h)\alpha)(\mathcal H_c)=\mathcal H_c\).
\end{enumerate}
For any one such \(h_0\), all aligning conjugators are exactly
\[
 \{h\in G:h\alpha(\mathcal H_c)h^{-1}=\mathcal H_c\}
        =\mathcal H_c h_0.
\tag{E.2.30c}
\]
For a peripheral-preserving automorphism \(\beta\), the condition
\(\beta(\mathcal H_c)=\mathcal H_c\) is equivalent to the two
memberships \(\beta(x),\beta(y)\in\mathcal H_c\).
Moreover,
\[
 \mathcal H_{\mathrm{int}}\subsetneq\mathcal H_c,
 \qquad
 x^{\widehat\Z},\ y^{\widehat\Z},\ z^{\widehat\Z}
       \subset\mathcal H_c.
\tag{E.2.30d}
\]
Applied to the actual Galois automorphisms in E.2.22, this is
an exact description of the proper-lattice premise. It does
not establish the affine boundary-extension vanishing of
E.2.28 or a discrete representative for the actual Galois action.
\end{result}
\begin{proof}
We retain the original cover lattices throughout and construct
the single aligning element from the actual quotient maps.

\paragraph{1. Peripheral filling and the embedded quotients.}
The group \(L_N\) is generated in \(K_N\) by its commutators
and its actual oriented cusp loops. If \(d_t\) is the order of
\(t\) in \(Q_N\), for \(t=x,y,z\), the latter are represented
by the words \(f t^{d_t}f^{-1}\), with \(f\) representing the
cosets of \(\langle t\rangle\) in \(Q_N\). Different path
choices change these words by conjugation inside \(K_N\),
which does not change the kernel. The deck group permutes the
cusps, so \(L_N\) is normal in \(F\). There is an exact sequence
\[
 1\longrightarrow A_N^c\longrightarrow\Gamma_N^c
       \longrightarrow Q_N\longrightarrow1.
\tag{E.2.30e}
\]
Here \(A_N^c\) is finite-rank free abelian. The group
\(\Gamma_N^c\) is residually finite: \(Q_N\) separates an
element outside \(A_N^c\), and \(\Gamma_N^c/mA_N^c\)
separates any nonzero element of \(A_N^c\), for a suitable
positive integer \(m\). The latter is a finite quotient,
since \(mA_N^c\) is normal and has finite index.
Consequently
\[
 F\cap J_N=L_N,\quad
 \Gamma_N^c\hookrightarrow P_N^c=\widehat{\Gamma_N^c},\quad
 N/J_N=\widehat A_N^c=A_N^c\otimes\widehat\Z,\quad
 \Gamma_N^c\cap\widehat A_N^c=A_N^c.
\tag{E.2.30f}
\]
The completion identity follows because a finite quotient of
\(F\) kills \(L_N\) exactly when its extension to \(G\)
kills \(J_N\). The quotients by \(mA_N^c\) show also that
the induced topology on \(A_N^c\) is its full profinite
topology. No torsion-freeness of \(\Gamma_N^c\) is required;
its peripheral images have finite order.

For \(M\subset N\), the inclusion \(K_M\to K_N\) sends
commutators to commutators and every cusp loop to a positive
integral power of a cusp loop. Hence \(L_M\subset L_N\) and
\(J_M\subset J_N\). The resulting maps
\(\Gamma_M^c\to\Gamma_N^c\) and \(P_M^c\to P_N^c\)
are surjective. Since \(J_N\subset N\), compatible coordinates
in the groups \(P_N^c\) define compatible nonempty compact
cosets in \(G\) with singleton intersection. Therefore
\[
 G\simeq\varprojlim_NP_N^c,\qquad
 \mathcal H_c\simeq\varprojlim_N\Gamma_N^c.
\tag{E.2.30g}
\]
Every projection \(\mathcal H_c\to\Gamma_N^c\) is onto,
because the included group \(F\) already projects onto it.
Giving the discrete quotients their discrete topologies makes
the second limit the completion of \(F\) for the neighborhoods
\(L_N\); the embedding in \(G\) is the one just constructed.

\paragraph{2. The original integral homology of every cover.}
The actual filling map for an arbitrary open \(U\subset G\)
induces a continuous surjection
\[
 \operatorname{ab}_U^c:U\longrightarrow
 \widehat\Lambda_U^c=\Lambda_U^c\otimes\widehat\Z
                  =\prod_\ell H_1(C_U,\Z_\ell).
\tag{E.2.30h}
\]
The last equality uses the actual integral Betti--\'{e}tale
comparison, or equivalently the completion of the discrete
surface homology lattice. For every \(U\),
\[
 \operatorname{ab}_U^c(\mathcal H_c\cap U)=\Lambda_U^c.
\tag{E.2.30i}
\]
Indeed, choose normal \(N\subset\operatorname{core}_G(U)\).
The map in (E.2.30h) kills \(L_N\), hence its closure \(J_N\).
For \(w\in\mathcal H_c\cap U\), choose \(f\in F\) with
\(wf^{-1}\in J_N\). Since \(J_N\subset U\), one has
\(f\in F\cap U\), and the proper homology classes of \(f\)
and \(w\) are equal. Conversely, \(F\cap U\) surjects onto
the original lattice \(\Lambda_U^c\).

Inner conjugation preserves these actual lattices between the
corresponding covers. For \(g\in G\), choose
\(f\in F\cap gU\) and write \(g=fu\), with \(u\in U\).
Then \(gUg^{-1}=fUf^{-1}\), and conjugation by \(u\) is
trivial on proper abelianization. The induced map is the
discrete change-of-basepoint map given by \(f\).

\paragraph{3. The zero invariant space and the self-normalizer.}
The affine regular cover retracts onto the graph with vertices
\(Q_N\) and oriented edges \((q,x),(q,y)\). Its invariant
integral edge cycles are exactly the span of
\[
 e_x=\sum_{q\in Q_N}(q,x),\qquad
 e_y=\sum_{q\in Q_N}(q,y).
\]
Both are boundary cycles: sum the oriented loops over all
cusps above \(x\), respectively above \(y\). Thus boundary
invariants surject onto affine homology invariants over \(\Q\).
The rational invariant functor of the finite group \(Q_N\)
is exact by averaging. Applying it to the actual boundary,
affine, and proper homology sequence gives
\[
 (A_N^c\otimes\Q)^{Q_N}=0.
\tag{E.2.30j}
\]
This also proves the quotient-curve equality with
\(H_1(C_N/Q_N,\Q)=H_1(\mathbf P^1,\Q)=0\) through
the actual filling maps.

Stack the integer matrices \(q-1\) on \(A_N^c\). Their
joint kernel over \(\Q\) is zero. Smith normal form then
shows that their joint kernel over every \(\Q_\ell\) and
over \(\widehat\Z\) is zero: multiplication by a nonzero
integer is injective on \(\widehat\Z\). In particular
\[
 (\widehat A_N^c)^{Q_N}=0.
\tag{E.2.30k}
\]
We apply the exact normalizer calculation of E.2.28 to
(E.2.30e). Write \(b=\gamma a\), with
\(\gamma\in\Gamma_N^c\) and \(a\in\widehat A_N^c\).
The condition that \(b\) normalize \(\Gamma_N^c\) is
\((q-1)a\in A_N^c\) for every \(q\in Q_N\). Smith normal
form produces an integral vector \(a_0\in A_N^c\) with
the same differences: an integral linear system with a
solution over \(\widehat\Z\) has an integral solution.
Then \(a-a_0\) is invariant, hence zero by (E.2.30k).
It follows that
\[
 \operatorname{Norm}_{P_N^c}(\Gamma_N^c)=\Gamma_N^c.
\tag{E.2.30l}
\]
If \(g\) normalizes \(\mathcal H_c\), its coordinate at
every \(N\) normalizes the surjective image \(\Gamma_N^c\),
so belongs to that image by (E.2.30l). Thus \(g\in\mathcal H_c\).
The reverse inclusion holds for every subgroup, proving
\[
 \operatorname{Norm}_G(\mathcal H_c)=\mathcal H_c.
\tag{E.2.30m}
\]

\paragraph{4. Peripheral automorphisms and compatible conjugators.}
Use the characteristic cofinal chain obtained by intersecting
the kernels of all homomorphisms to symmetric groups on at
most \(n\) letters, as in E.2.28. Every \(\alpha\) fixes
its members \(N\). The subgroup \(J_N\) is intrinsically
the closed normal subgroup generated by the commutators of
\(N\) and the cusp powers belonging to \(N\). Peripheral
preservation gives \(\alpha(J_N)=J_N\): the induced finite
quotient automorphism preserves the relevant cusp orders,
and a unit power generates the same closed cusp subgroup.
Using \(\alpha^{-1}\) proves equality. Hence \(\alpha\)
induces an actual automorphism \(\alpha_N\) of \(P_N^c\).

Assume condition 2. The abstract group
\(V_N=\widehat A_N^c/A_N^c\) is uniquely divisible.
Surjectivity of multiplication by \(m>0\) follows by
subtracting an integral residue modulo \(m\). Injectivity
follows because an integer divisible by \(m\) in every
\(\Z_\ell\) is divisible by \(m\) in \(\Z\).
Thus \(V_N\) is a \(\Q\)-vector space.
The two groups \(\Gamma_N^c/A_N^c\) and
\(\alpha_N(\Gamma_N^c)/A_N^c\) are complements in the
abstract extension \(P_N^c/A_N^c\) of \(Q_N\) by \(V_N\).
Their difference is a cocycle \(c:Q_N\to V_N\). With
\[
 v=\frac{1}{|Q_N|}\sum_{q\in Q_N}c(q),
 \qquad c(q)=v-qv,
\]
conjugation by any lift of \(-v\) carries the second
complement to the first. Before quotienting, both subgroups
contain \(A_N^c\), so this is an actual local aligning
conjugator in \(P_N^c\).

Define its full solution set by
\[
 S_N^c=\{u\in P_N^c:
        \operatorname{ad}(u)\alpha_N(\Gamma_N^c)=\Gamma_N^c\}.
\tag{E.2.30n}
\]
By (E.2.30l) this nonempty set is \(\Gamma_N^c u_N\)
for any one solution \(u_N\), the orbit for left
multiplication by \(\Gamma_N^c\). The transition maps
\(S_M^c\to S_N^c\) are onto. Indeed, choose any solution
above \(M\), compare its image with the desired solution
above \(N\), and lift their ratio through
\(\Gamma_M^c\twoheadrightarrow\Gamma_N^c\).
Selecting a compatible sequence along the countable cofinal
chain gives a single \(h\in G\) by (E.2.30g). Its coordinates
align every \(\Gamma_N^c\), so
\((\operatorname{ad}(h)\alpha)(\mathcal H_c)=\mathcal H_c\).
This proves \(2\Rightarrow3\) without applying compactness
to any nonclosed solution set. Formula (E.2.30i) and the
actual inner-conjugation homology maps prove \(3\Rightarrow1\);
the implication \(1\Rightarrow2\) is immediate.

Two global solutions \(h,h_0\) differ by an element
\(hh_0^{-1}\) normalizing \(\mathcal H_c\). Formula
(E.2.30m) gives exactly (E.2.30c), including uniqueness
of the aligning coset in \(\mathcal H_c\backslash G\).

\paragraph{5. The generator criterion and the strict inclusion.}
If \(\beta(x),\beta(y)\in\mathcal H_c\), then
\(\beta_N(\Gamma_N^c)\subset\Gamma_N^c\) at every
characteristic level. Its restriction to \(A_N^c\) is
an integral matrix whose completion is an automorphism.
Its determinant is therefore an integer invertible at every
prime, hence \(\pm1\). It maps \(A_N^c\) onto itself.
The induced map on \(Q_N\) is onto as well, and the exact
extension (E.2.30e) gives
\(\beta_N(\Gamma_N^c)=\Gamma_N^c\). In rank zero this
uses the unique automorphism of the zero lattice. Taking
the embedded inverse limit proves preservation of
\(\mathcal H_c\). The converse follows because
\(x,y\in\mathcal H_c\).

The affine commutator kernel \(D_N\) is contained in
\(J_N\); its actual quotient maps give
\(\mathcal H_{\mathrm{int}}\subset\mathcal H_c\).
At each proper level the image of a peripheral generator
has finite order dividing its inertia order in \(Q_N\).
Every profinite power therefore has the same coordinate
as an integer power at that level. This proves all three
peripheral inclusions in (E.2.30d).

For strictness choose \(a\in\widehat\Z^{\times}\) with
\(a_2=-1\) and \(a_p=1\) at every odd prime. It is not
an integer: its \(3\)-adic coordinate would force that
integer to be \(1\), contradicting its \(2\)-adic
coordinate. Then \(x^a\in\mathcal H_c\). At the affine
level \(N=G\), however, the image of
\(\mathcal H_{\mathrm{int}}\) in
\(G^{\mathrm{ab}}=\widehat\Z^2\) lies in \(\Z^2\).
Thus \(x^a\notin\mathcal H_{\mathrm{int}}\).

This element does not supply a Galois automorphism or an
integral transporter. For the actual \(\sigma\in\Gamma_{F_{\mathrm{CM}}}\)
of E.2.22 the cyclotomic character is one, so its cusp
generators are transported to conjugates of the same
generators. The remaining proper-to-affine assertion is
the inclusion, inside this actual Galois image, of the
conjugacy-stabilizer condition for \(\mathcal H_c\) in
the corresponding condition for \(\mathcal H_{\mathrm{int}}\).
The strict subgroup inclusion just proved establishes
neither that assertion nor its failure. The exact
off-diagonal boundary-extension block in E.2.28 and the
discrete-representative requirement remain to be resolved.
\end{proof}
